\section{Table and figure extraction}\label{section:background-table-and-figure-extraction}

Figures, and especially tables, are the parts of a paper that are the most difficult for a layout parser to recover. However, the central quantitative evidence is presented in such formats in histopathology papers. Table extraction is commonly divided into two separate tasks: table detection, and table structure recognition. Table detection involves locating the tables and their boundaries, whereas table structure recognition involves the reconstruction of rows, columns, and cells in the table. Recent work on these tasks increasingly uses image-based transformer detectors~\parencite{}. Figure extraction and the association of captions to figures follow a similar line, with systems designed to detect figures, extract their captions, and associate them correctly~\parencite{}.

The thesis follows the logic in recent work, separating the steps of detection, structure recognition, and association. The document-extraction stage combines a layout-based detector (Docling) with an image-based table detector (Microsoft's Table Transformer) in a hybrid detector, and adds footnote-expansion, geometric caption-matching, and decorative artifact filtering to the system, as described in Section~\ref{section:Document_Extraction}. The results in Chapter~\ref{chapter:Results} show that off-the-shelf Docling already localizes tables and figures well; the improvements over the off-the-shelf Docling baseline mostly come from footnote capture for tables, and from caption attachment and decorative-icon filtering for figures.